La recherche présentée se concentre sur l'emploi de l'intelligence artificielle de type LLM pour la génération et le montage de clips vidéo dans le contexte du marketing d’influence. Ce type de marketing paie des influenceurs pour promouvoir des produits, et leur salaire est proportionnel au nombre de vues mensuelles des vidéos qu’ils produisent.
Le projet se divise en deux parties, mais cette recherche ne portera que sur la deuxième. La première concerne le développement d’un programme autonome pour remplacer le travail d'un monteur de vidéos humain par une IA appelée “monteur”, en utilisant des intelligences artificielles de type LLM et vision. La deuxième partie consiste au développement d’un programme autonome pour remplacer le travail d’un “directeur” artistique en utilisant des intelligences artificielles de type LLM.
Mener cette recherche à terme passera par la création d’une intelligence artificielle capable de prendre des décisions artistiques sur la sélection des moments forts des vidéos pour créer la base des clips et les coupures à réaliser pour augmenter le nombre de vues. Cela impliquera également de découvrir le type de montage qui maximise le nombre de vues en instaurant un système de tests A/B sur des types de montages préalablement créés et catégorisés par l’équipe de recherche, avec un retour sur les vidéos générés, et en intégrant automatiquement cet apprentissage à l’intelligence artificielle du “directeur”.
Les retombées attendues incluent, d'une part, le développement d'une méthode innovante permettant de convertir des vidéos en un format textuel structuré, fournissant toutes les informations essentielles pour permettre leur traitement par un modèle linguistique purement textuel. Cette approche ouvre des perspectives pour exploiter des modèles textuels dans des contextes multi-modaux.
D'autre part, cette recherche propose une nouvelle technique d’apprentissage où l’intelligence artificielle "directeur" s’améliore de manière continue en intégrant les données issues de la performance des vidéos qu’elle produit. Ce cadre méthodologique, basé sur l’adaptation en temps réel, pourra être réutilisé par la communauté scientifique pour concevoir des systèmes autonomes capables d’apprendre directement de leurs interactions avec leur environnement.
